\documentclass[]{article}
\usepackage{tikz,pgfplots}
\usepackage{geometry}
\usepackage{float}

\begin{document}
%\section{Drawing lines and Scaling picture}
%\begin{figure}
%\begin{tikzpicture}
%\draw(0,0)--(1,2)--(2,3)--(1,0);
%\draw[help lines](0,0) grid(2,3);
%\end{tikzpicture}
%\caption{First figure}
%\end{figure}

%\begin{figure}
%\begin{tikzpicture}[scale=3]
%\draw (0,0)--(1,2);
%\end{tikzpicture}
%\end{figure}

%\begin{figure}
%\begin{tikzpicture}[xscale=3]
%\draw (0,0)--(2,0);
%\end{tikzpicture}
%\end{figure}

%\begin{figure}[H]
%\begin{tikzpicture}[yscale=3]
%\draw (0,0)--(2,0);
%\end{tikzpicture}
%\end{figure}

%\begin{figure}
%\begin{tikzpicture}[xscale=3,yscale=0.7]
%\draw (0,0)--(2,0);
%\end{tikzpicture}
%\end{figure}

%\begin{figure}
%\begin{tikzpicture}[xscale=3]
%\draw[->] (0,0)--(2,0);
%\draw[<-] (1,1)--(2,1);
%\draw[| ->] (2,2)--(4,2);
%\end{tikzpicture}
%\end{figure}

%\begin{figure}[H]
%\begin{tikzpicture}
%\draw[<->](0,2)--(0,0)--(3,0);
%\end{tikzpicture}
%\end{figure}

%everything we can do with draw
%\begin{figure}
%\begin{tikzpicture}
%\draw[<->,thick,gray](0,0)--(2,1);
%\draw[line width=0.3cm](2,1)--(2,3);
%\end{tikzpicture}
%\end{figure}
%\subsection{changing thickness}
%\begin{tikzpicture}
%\draw[<->](0,2)--(0,0)--(4,0);
%\draw[very thick](0,1)--(3,0);
%\draw[thick](0,0)--(2,3);
%\draw[help lines](0,1)--(1,1)--(1,0);
%\end{tikzpicture}

%\subsection{curve}
%\begin{figure}
%\begin{tikzpicture}
%\draw[thick,gray] (0,0) rectangle (4,2);
%\draw[blue!30] (2,2) circle [radius=2cm];
%\draw[fill,pink!80](2,2) arc[radius=1cm,start angle=90,end angle=275];
%\draw[<->,rounded corners] (2,2)--(2,0)--(4,0);
%\draw[](5,5) to[out=90,in=180] (8,9);
% \draw [<->,thick, cyan,fill] (0,0) to [out=90,in=180] (1,1)
% to [out=0,in=180] (2.5,0) to [out=0,in=-135] (4,1) ;
%\end{tikzpicture}
%\end{figure}

%\subsection{Plots}
%
%\begin{figure}[H]
%\begin{tikzpicture}
%\draw[very thick,->] (0,0)--(10,0);%x line
%\draw[very thick,->] (0,0)--(0,10);%y line
%\draw[step=1.5cm,ultra thin] (0,0)  grid (9.9,9.9);
%\draw[thick,blue!40,domain=0:2.5] plot(\x,{0.25+\x+\x*\x});
%\draw[very thick,pink!80,domain=0:9] plot(\x,{\x/2});
%\draw[very thick,cyan!90,domain=0:2*pi] plot(\x,{sin(\x r)});
%\draw[very thick,red!40,domain=0:2*pi] plot(\x,{cos(\x r)});
%\draw[very thick,green!30,domain=0:2] plot (\x,{exp(\x)});
%\draw[very thick,yellow!80,domain=1:9] plot(\x,{ln(\x)});%in ln function domain should start from 1 cause 0 is invalid
%\end{tikzpicture}
%\end{figure}

%\section{Filling}
%\begin{figure}
%\begin{tikzpicture}
%\draw[thick,fill=pink!40](0,0) rectangle (4,2);
%\path[thick,fill=pink!80](5,0) rectangle(10,2); %if we don't want to see the outline then we should use path
%\end{tikzpicture}
%\end{figure}

%\section{Node and labeling}
%\begin{figure}[H]
%\begin{tikzpicture}
%\draw[ultra thick,->] (0,0)--(10,0);
%\draw[ultra thick,->](0,0)--(0,10);
%\draw[ultra thin,step=1.5cm] (0,0) grid (9.9,9.9);
%\node[below right] at (10,0) {$X$};
%\node[above left] at(0,10) {$Y$};
%\node[below right] at(8,0){\mbox{(step=1.5cm)}};
%\draw[fill=pink!30] (5,5) circle [radius=1.5cm];
%\node[below,cyan]at (5,5) {\mbox{(Pink circle)}};
%\end{tikzpicture}
%\end{figure}

\section{Example}
%\begin{figure}[H]
%\begin{tikzpicture}
%\draw[ultra thick,red](0,0)--(5,0);
%\draw[ultra thick,green](5,0)--(10,0);
%\draw[ultra thick](0,-0.1)--(0,0.1);
%\draw[ultra thick](5,-0.1)--(5,0.1);
%\draw[ultra thick](10,-0.1)--(10,0.1);
%\node[below] at (5,0) {$a=b=1/2$};
%\node[below] at (0,0) {0};
%\node[below] at (10,0){1};
%\end{tikzpicture}
%\end{figure}

\begin{figure}[H]
\begin{tikzpicture}
\draw[thick,red](0,0)--(5,0);
\draw[thick,red](0,-1)--(5,-1);
\draw[][thick,green](5,-1)--(10,-1);

\end{tikzpicture}
\end{figure}
\end{document}
